\documentclass[a4paper,11pt]{article}
\usepackage[utf8]{inputenc}
\usepackage[T1]{fontenc}
\usepackage[english]{babel}
\usepackage{csquotes} % Requis pour biblatex
\usepackage[backend=biber,style=apa]{biblatex}
\addbibresource{PathPlanning.bib}

\usepackage{amsmath, amssymb}
\usepackage{graphicx}
\usepackage{listings}
\usepackage{xcolor}
\usepackage{caption}
\usepackage{hyperref}

% Code source style
\lstdefinestyle{mystyle}{
    backgroundcolor=\color{gray!10},   
    commentstyle=\color{green!50!black},
    keywordstyle=\color{blue},
    numberstyle=\tiny\color{gray},
    stringstyle=\color{orange},
    basicstyle=\ttfamily\small,
    breaklines=true,
    frame=single,
    numbers=left,
    numbersep=5pt
}
\lstset{style=mystyle}

\title{Exemple d'article scientifique en \LaTeX}
\author{Ton Nom}
\date{\today}

\begin{document}

\maketitle

\begin{abstract}
    Cet article montre un exemple minimal d’utilisation de \LaTeX{} avec Zotero, du code, des mathématiques et des figures.
\end{abstract}

\section{Introduction}

Voici une citation tirée de Zotero \cite{wuResearchProcessPath2024}.
On peut écrire des maths comme : \( E = mc^2 \) ou
\[
    \int_0^1 x^2\,dx = \frac{1}{3}
\]

\section{Code source}

Un extrait de code Python :

\begin{lstlisting}[language=Python, caption=Exemple de code Python]
def fib(n):
    if n <= 1:
        return n
    return fib(n-1) + fib(n-2)
\end{lstlisting}

\section{Figure}

Voici une image (placer `exemple.png` dans le dossier `images/`) :

% \begin{figure}[h]
%     \centering
%     \includegraphics[width=0.5\textwidth]{images/exemple.png}
%     \caption{Exemple d'image insérée dans LaTeX}
% \end{figure}

\section{Conclusion}

Ce modèle vous permet d’écrire efficacement un article scientifique court en LaTeX.

\printbibliography

\end{document}

% \documentclass{article}
% \usepackage[utf8]{inputenc}
% \usepackage[backend=biber,style=apa]{biblatex}
% \addbibresource{Path\ planning.bib}

% \title{Test demo}
% \author{Pauline Prost--Waterson}
% \begin{document}
% \maketitle
% \section{A* Algorithm}
% So I started with A* algorithm. Complexity is $O(b^d)$ where $b$ is the branching factor and $d$ is the depth of the solution.
% \section{RRT Algorithm}
% Complexity is $O(n \log n)$ where $n$ is the number of nodes in the tree. I used binary search tree to find the nearest node.
% Description of the algorithm + article source : \cite{wuResearchProcessPath2024}.
% \section{RRT* Algorithm}
% jsp quoi écrire de plus

% \printbibliography
% \end{document}